%---------------------------------------------------------------------
% Term glossary and the recurring mistakes.
%
% The glossary matters more in this subject than in a mathematics
% course: the paper expects the English technical term, but the
% explanation has to be given in Urdu, and most marks are lost in the
% gap between the two.
%---------------------------------------------------------------------

\section*{اصطلاحات کی فہرست}
\label{sec:lughat}
\markboth{اصطلاحات کی فہرست}{}

پرچے میں اصطلاح انگریزی ہی میں لکھیے، اور ساتھ ایک سطری اردو وضاحت دیجیے۔ صرف
ترجمہ لکھ دینے سے نمبر پورے نہیں ملتے۔

\subsection*{ہارڈ ویئر اور ساخت}

\begin{center}
\begin{tabular}{@{}p{4.0cm} p{10.4cm}@{}}
\toprule
\textenglish{Hardware} & کمپیوٹر کے وہ اجزا جنہیں چھوا جا سکتا ہے \\
\textenglish{Software} & ہدایات کا مجموعہ جو کمپیوٹر کو کام بتاتا ہے \\
\textenglish{CPU} & \textenglish{Central Processing Unit} --- کمپیوٹر کا دماغ \\
\textenglish{ALU} & \textenglish{Arithmetic and Logic Unit} --- حسابی و منطقی عمل \\
\textenglish{CU} & \textenglish{Control Unit} --- ہدایات کی ترتیب اور نگرانی \\
\textenglish{RAM} & عارضی میموری؛ بجلی جاتے ہی خالی (\textenglish{volatile}) \\
\textenglish{ROM} & مستقل میموری؛ صرف پڑھی جا سکتی ہے \\
\textenglish{Port} & بیرونی آلہ جوڑنے کا نقطہ (\textenglish{USB}، \textenglish{HDMI}) \\
\textenglish{Bus} & اندرونی راستہ جس پر ڈیٹا ایک حصے سے دوسرے تک جاتا ہے \\
\textenglish{Microprocessor} & ایک \textenglish{chip} پر بنا مکمل \textenglish{CPU} \\
\bottomrule
\end{tabular}
\end{center}

\subsection*{اِن پُٹ اور آؤٹ پُٹ}

\begin{center}
\begin{tabular}{@{}p{4.0cm} p{10.4cm}@{}}
\toprule
\textenglish{OCR} & \textenglish{Optical Character Recognition} --- حروف پہچان کر متن بناتا ہے \\
\textenglish{OMR} & \textenglish{Optical Mark Recognition} --- خانوں میں لگے نشان پڑھتا ہے \\
\textenglish{MICR} & \textenglish{Magnetic Ink Character Recognition} --- بینک چیک کے حروف \\
\textenglish{BCR} & \textenglish{Bar Code Reader} --- اشیا پر چھپی لکیریں پڑھتا ہے \\
\textenglish{Impact Printer} & کاغذ پر ضرب لگاتا ہے (\textenglish{Dot Matrix}) \\
\textenglish{Non-Impact Printer} & ضرب نہیں لگاتا (\textenglish{Inkjet}، \textenglish{Laser}) \\
\textenglish{Plotter} & قلم سے بڑے نقشے اور خاکے بناتا ہے \\
\textenglish{POS Terminal} & خریداری کے مقام پر قیمت، رسید اور \textenglish{inventory} \\
\textenglish{ATM} & \textenglish{Automated Teller Machine} --- بینک کی خودکار مشین \\
\textenglish{Resolution} & سکرین یا تصویر کی وضاحت، \textenglish{pixels} میں \\
\bottomrule
\end{tabular}
\end{center}

\subsection*{سافٹ ویئر، آپریٹنگ سسٹم اور زبانیں}

\begin{center}
\begin{tabular}{@{}p{4.0cm} p{10.4cm}@{}}
\toprule
\textenglish{Operating System} & ہارڈ ویئر اور صارف کے درمیان واسطہ؛ وسائل کا منتظم \\
\textenglish{GUI} & \textenglish{Graphical User Interface} --- شبیہوں کے ذریعے رابطہ \\
\textenglish{Utility Program} & نظام کی دیکھ بھال کا سافٹ ویئر (\textenglish{antivirus}) \\
\textenglish{Process Management} & کئی پروگراموں کو باری باری \textenglish{CPU} دینا \\
\textenglish{Compiler} & پورا پروگرام ایک ساتھ ترجمہ؛ \textenglish{object code} بناتا ہے \\
\textenglish{Interpreter} & ایک وقت میں ایک سطر کا ترجمہ \\
\textenglish{Assembler} & \textenglish{assembly language} کو مشینی زبان میں بدلتا ہے \\
\textenglish{Linker} & کئی \textenglish{object files} جوڑ کر ایک قابلِ عمل پروگرام \\
\bottomrule
\end{tabular}
\end{center}

\subsection*{نیٹ ورکنگ اور ملٹی میڈیا}

\begin{center}
\begin{tabular}{@{}p{4.0cm} p{10.4cm}@{}}
\toprule
\textenglish{LAN / MAN / WAN} & مقامی / شہری / وسیع علاقے کا نیٹ ورک \\
\textenglish{Topology} & نیٹ ورک میں کمپیوٹروں کی ترتیب کا نقشہ \\
\textenglish{Protocol} & رابطے کے قواعد، جن پر دونوں فریق متفق ہوں \\
\textenglish{OSI Model} & نیٹ ورکنگ کا سات پرتوں والا معیاری خاکہ \\
\textenglish{Switching} & \textenglish{LAN} کے اندر ڈیٹا صحیح آلے تک پہنچانا \\
\textenglish{Routing} & مختلف نیٹ ورکوں کے درمیان بہترین راستے کا انتخاب \\
\textenglish{NIC} & \textenglish{Network Interface Card} --- نیٹ ورک سے جوڑنے والا کارڈ \\
\textenglish{Multimedia} & متن، تصویر، آواز، ویڈیو اور \textenglish{animation} کا امتزاج \\
\textenglish{Animation} & مصنوعی حرکت، ویڈیو کے برعکس \\
\bottomrule
\end{tabular}
\end{center}

\medskip
\begin{nukta}[\textenglish{OSI} کی سات پرتیں --- ایک سوال کے لیے کافی]
نیچے سے اوپر: \textenglish{Physical}، \textenglish{Data Link}،
\textenglish{Network}، \textenglish{Transport}، \textenglish{Session},
\textenglish{Presentation}، \textenglish{Application}۔ ترتیب یاد رکھنے کے لیے
انگریزی جملہ: \textenglish{\emph{Please Do Not Throw Sausage Pizza Away}}۔
\end{nukta}

\clearpage

%---------------------------------------------------------------------
\section*{آٹھ عام غلطیاں}
\markboth{آٹھ عام غلطیاں}{}

\begin{enumerate}\setlength{\itemsep}{5pt}
  \item \textbf{صرف تعریف لکھ کر رک جانا۔} اس مضمون کا ہر سوال تعریف کے ساتھ
        \emph{مثال} مانگتا ہے۔ دو مثالوں کے بغیر جواب ادھورا ہے۔
  \item \textbf{''فرق واضح کیجیے'' کا جواب پیراگراف میں لکھنا۔} جدول بنائیے۔
        ممتحن کو نکات فوراً نظر آتے ہیں اور نمبر پورے ملتے ہیں۔
  \item \textbf{\textenglish{OCR} اور \textenglish{OMR} کو خلط ملط کرنا۔}
        \textenglish{OCR} \emph{حروف} پڑھتا ہے، \textenglish{OMR} صرف
        \emph{نشان}۔
  \item \textbf{تاریخ اور نسل کو ایک سمجھنا۔} نسل کی بنیاد ہمیشہ اُس دور کا
        \emph{الیکٹرانک پرزہ} ہوتا ہے: \textenglish{vacuum tube}،
        \textenglish{transistor}، \textenglish{IC}، \textenglish{microprocessor}۔
  \item \textbf{\textenglish{Volatile} اور \textenglish{Non-Volatile} الٹ
        دینا۔} \textenglish{RAM} بجلی جاتے ہی خالی؛ \textenglish{ROM} اور ہارڈ
        ڈسک محفوظ۔
  \item \textbf{\textenglish{System} اور \textenglish{Application software} کی
        مثالیں بدل دینا۔} \textenglish{Windows} سسٹم ہے،
        \textenglish{MS Word} ایپلیکیشن۔
  \item \textbf{\textenglish{Topology} میں صرف نام گنوا دینا۔} ہر ترتیب کے ساتھ
        ایک خوبی اور ایک خامی لکھیے، اور ہو سکے تو خاکہ بنائیے۔
  \item \textbf{اصطلاح کا اردو ترجمہ لکھ دینا۔} تکنیکی اصطلاح انگریزی ہی میں
        لکھیے (\textenglish{Full Duplex})، اور وضاحت اردو میں دیجیے۔
\end{enumerate}

\vfill
\begin{center}
\begin{tcolorbox}[enhanced, colback=bmtint, colframe=bmblue, boxrule=0.5pt,
                  arc=3pt, width=0.94\linewidth,
                  left=11pt, right=11pt, top=9pt, bottom=9pt]
\small
\textbf{\color{bmblue}اگر وقت کم ہو۔}\ \
یہ پانچ موضوعات سب سے زیادہ نفع بخش ہیں:
\textenglish{(1)} \textenglish{CPU}، \textenglish{ALU}، \textenglish{CU} اور
میموری کی اقسام؛
\textenglish{(2)} \textenglish{OCR / OMR / MICR / BCR} کا تقابل؛
\textenglish{(3)} پرنٹر کی اقسام اور \textenglish{impact} بمقابلہ
\textenglish{non-impact}؛
\textenglish{(4)} \textenglish{System} بمقابلہ \textenglish{Application
software} اور آپریٹنگ سسٹم کے افعال؛
\textenglish{(5)} \textenglish{LAN/WAN}، \textenglish{topologies} اور
\textenglish{Simplex/Duplex}۔

\medskip
\textbf{\color{bmblue}ماخذ۔}\ \
\textenglish{\emph{Basics of Information and Communication Technology}}،
کورس کوڈ \textenglish{1431/5403}، یونٹ \textenglish{1--9}، شعبۂ کمپیوٹر سائنس،
علامہ اقبال اوپن یونیورسٹی، اسلام آباد (پہلا ایڈیشن \textenglish{2013})۔ تمام
سوالات اسی کتاب کے \textenglish{Self-Assessment Questions} سے لیے گئے ہیں۔
\end{tcolorbox}
\end{center}
